\documentclass[12pt, a4paper]{article}

\usepackage[utf8]{inputenc}
\usepackage{geometry}
\geometry{a4paper, margin=1in}
\usepackage{graphicx}
\usepackage{hyperref}
\usepackage{booktabs}
\usepackage{tabularx}
\usepackage{float}

\title{
    \vspace{2cm}
    \textbf{Software Engineering Project Report}\\
    \vspace{1cm}
    \Large \textbf{Lost \& Found Management System (LFMS)}
    \vspace{2cm}
}
\author{
    \textbf{Group 9}\\
    SE VLabs Institute
}
\date{\today}

\begin{document}

\maketitle
\newpage

\tableofcontents
\newpage

\section{Introduction}
\subsection{Purpose}
The purpose of this Software Requirements Specification (SRS) and project report is to describe in detail the development and implementation of the Lost \& Found Management System (LFMS). This document specifies the functional and non-functional requirements, architectural concepts, and the technical implementation details of the system.

\subsection{Scope}
The LFMS is a web-based application designed to manage lost and found items within the institute campus. Currently, lost-and-found activities are managed through manual verbal communication, notice boards, and email, which leads to inefficiencies, mismanagement, and limited transparency. 
The developed system provides:
\begin{itemize}
    \item A centralized digital platform
    \item User-based reporting of lost and found items
    \item Administrative verification before publication
    \item A secure claim request handling and approval system
\end{itemize}

\subsection{Definitions \& Acronyms}
\begin{itemize}
    \item \textbf{LFMS}: Lost \& Found Management System
    \item \textbf{LAN}: Local Area Network
    \item \textbf{DSW}: Dean of Student Welfare
    \item \textbf{SRS}: Software Requirements Specification
    \item \textbf{Admin}: Authorized officer responsible for verification
\end{itemize}

\section{Overall Description}
\subsection{Product Perspective}
The LFMS is a standalone, web-based application that replaces the current manual and email-based reporting system. It follows a modern Client-Server architecture where users interact via an HTML5 web browser, and the backend server processes requests, handling authentication and database interactions securely.

\subsection{User Classes and Characteristics}
\begin{itemize}
    \item \textbf{Student / Registered User}: Can register, log in, post lost items, submit claim requests, and view their dashboard.
    \item \textbf{Administrator (DSW / Authorized Officer)}: Verifies lost and found posts, approves or rejects claims, manages user accounts (e.g., reactivation), and has full system control privileges.
    \item \textbf{Guest User (Non-Registered)}: Can browse and search available items but cannot post or claim items.
\end{itemize}

\subsection{Operating Environment}
The system is deployed as a web application accessible within the institute LAN. It is fully compatible with modern browsers including Google Chrome, Mozilla Firefox, Safari, and Opera.

\section{System Architecture \& Software Engineering Concepts}
The project applies fundamental Software Engineering methodologies, utilizing standard UML modeling to capture requirements and system flow.

\subsection{Use Case Modeling}
The system functionality was captured using Use Case scenarios mapping actors (Student, Admin, Guest) to actions:
\begin{itemize}
    \item \textbf{UC1 - Register}: Allows new users to create accounts securely.
    \item \textbf{UC2 - Login}: Authenticates users, featuring an \textit{«extend»} relationship for handling security questions upon 3 consecutive failed attempts.
    \item \textbf{UC3 - Search Items}: Guest and Registered users can browse the directory.
    \item \textbf{UC4 - Post Lost Item}: Logged-in users report items they lost.
    \item \textbf{UC5 - Register Found Item}: Logged-in users report found items.
    \item \textbf{UC6 - Admin Verification}: Admin reviews and publishes/rejects found items (\textit{«include»} in the publication flow).
    \item \textbf{UC7 - Submit Claim}: Owner provides proof for a found item.
    \item \textbf{UC8 - Approve Claim}: Admin verifies proof and marks item as returned.
\end{itemize}

\subsection{State Transition Table}
The lifecycle of a "Found Item" is managed through the following state transitions:
\begin{table}[H]
    \centering
    \begin{tabularx}{\textwidth}{|l|l|X|}
        \hline
        \textbf{Initial State} & \textbf{Event} & \textbf{Next State} \\
        \hline
        None & User Reports Item & PENDING\_VERIFICATION \\
        \hline
        PENDING\_VERIFICATION & Admin Approves & PUBLISHED \\
        \hline
        PUBLISHED & User Submits Claim & CLAIM\_PENDING \\
        \hline
        CLAIM\_PENDING & Admin Approves Claim & RETURNED \\
        \hline
        CLAIM\_PENDING & Admin Rejects Claim & PUBLISHED \\
        \hline
    \end{tabularx}
    \caption{Found Item State Machine}
\end{table}

\subsection{Entity Relationship (ER) Model}
The database schema involves three primary entities with the following relationships:
\begin{itemize}
    \item \textbf{User (1) $\rightarrow$ Item (N)}: A user can report multiple lost or found items.
    \item \textbf{User (1) $\rightarrow$ Claim (N)}: A user can submit multiple claims for different found items.
    \item \textbf{Item (1) $\rightarrow$ Claim (N)}: A single found item can have multiple competing claims which are reviewed by the admin.
\end{itemize}
\begin{figure}[H]
    \centering
    \fbox{
        \begin{minipage}{0.8\textwidth}
            \centering
            \textbf{[User]} ---owns--- \textbf{[Item]}\\
            \textbf{[User]} ---submits--- \textbf{[Claim]}\\
            \textbf{[Item]} ---receives--- \textbf{[Claim]}
        \end{minipage}
    }
    \caption{Conceptual ER Model Overview}
\end{figure}

\subsection{System Architecture}
The system follows a \textbf{Model-View-Controller (MVC)} pattern implemented via the Next.js framework:
\begin{itemize}
    \item \textbf{Model}: Prisma Schemas (SQLite)
    \item \textbf{View}: React Components (Next.js App Router)
    \item \textbf{Controller}: Next.js API Routes and Server Actions
\end{itemize}

\section{Functional Requirements}
\begin{table}[H]
    \centering
    \begin{tabularx}{\textwidth}{|l|X|l|}
        \hline
        \textbf{ID} & \textbf{Requirement Description} & \textbf{Priority} \\
        \hline
        R1 & \textbf{New User Registration}: Unique user ID generation; hashed passwords. & High \\
        \hline
        R2 & \textbf{User Login}: Authentication with 3-attempt lockout and security question reset mechanism. & High \\
        \hline
        R3 & \textbf{Search Item}: Search by name, category, location, and date (Open to guests). & High \\
        \hline
        R4 & \textbf{Post Lost Item}: Logged-in users can report lost items. & High \\
        \hline
        R5 & \textbf{Post Found Item}: Users can report found items; requires Admin verification. & High \\
        \hline
        R6 & \textbf{Claim Item}: Users can submit proof to claim a found item. & High \\
        \hline
        R7 & \textbf{Admin Approval}: Admin reviews posts, approves/rejects claims. & High \\
        \hline
    \end{tabularx}
    \caption{Functional Requirements Mapping}
\end{table}

\section{Non-Functional Requirements}
\begin{itemize}
    \item \textbf{Performance}: The system easily supports 50+ concurrent users with sub-second response times using modern edge-ready server logic.
    \item \textbf{Security}: Sensitive data (passwords, security answers) are encrypted using \texttt{bcrypt}. Authentication is maintained via secure HTTP-only JSON Web Tokens (JWT).
    \item \textbf{Usability}: The frontend was designed using a minimalist and elegant monochromatic theme, focusing on typography and clear accessibility.
    \item \textbf{Maintainability}: The codebase utilizes a modular, component-based React architecture allowing for rapid future enhancements.
\end{itemize}

\section{Technical Implementation}
To fulfill the requirements, modern web development frameworks were utilized:
\begin{itemize}
    \item \textbf{Frontend}: HTML5, React, and Next.js utilizing a custom minimalist Vanilla CSS design system.
    \item \textbf{Backend}: Next.js App Router API Routes serving as the controller layer.
    \item \textbf{Database}: SQLite integrated with the Prisma Object-Relational Mapper (ORM) for type-safe database queries and automated migrations.
\end{itemize}

\section{Conclusion}
The Lost \& Found Management System was successfully developed and deployed as a structured, transparent, and secure solution. By replacing the existing manual and email-based processes with this centralized web portal, the SE VLabs Institute ensures better tracking, faster recovery, strict administrative control, and a vastly improved user experience for its students and staff.

\end{document}
